\documentclass{article}
\usepackage{lmodern}
\usepackage{amsmath}
\usepackage{listings}
\usepackage{fullpage}
\usepackage{parskip}
\usepackage{xcolor}
\lstset{
	basicstyle=\ttfamily,
	columns=fullflexible,
	frame=single,
	breaklines=true,
	postbreak=\mbox{\textcolor{red}{$\hookrightarrow$}\space},
}
\begin{document}
\title{Experiment `p\_6\_4\_sum\_without\_smallest\_array' Results}
\author{\today}
\date{}
\maketitle
\textbf{Experiment outcome: }SUCCESS
\\\textbf{Bad responses: }0
\\\textbf{Responses containing}~\texttt{assume}~\textbf{: }0
\\\textbf{Responses containing}~\texttt{expect}~\textbf{: }0
\\\textbf{Resolution attempts: }1
\\\textbf{Hard fails (resolution): }0
\\\textbf{Soft fails (resolution): }0
\\\textbf{Verification attempts: }1
\section*{Problem Specification}
\textbf{Problem name: }p\_6\_4\_sum\_without\_smallest\_array
\\\textbf{Natural language statement: }Write a method that computes the sum of an array of values, except for the smallest one, in a single loop. In the loop, update the sum and the smallest value. After the loop, return the difference.
\\\textbf{Method signature: }p\_6\_4\_sum\_without\_smallest\_array(inputs: array<int>) returns (sum: int)
\subsection*{Ensures}
\begin{itemize}
\item \texttt{sum == sum\_array(inputs[..]) - find\_smallest(inputs[..])}
\end{itemize}
\subsection*{Requires}
\begin{itemize}
\item \texttt{inputs.Length >= 1}
\end{itemize}
\subsection*{Functional Code Given}
\begin{lstlisting}
function find_smallest(inputs: seq<int>): (ret: int)
  requires |inputs| >= 1
  ensures forall x :: x in inputs ==> ret <= x 
{
  if |inputs| == 1 then inputs[0]
  else 
    assert 2 <= |inputs|;
    var rest_smallest := find_smallest(inputs[1..]);
    var ret: int := if inputs[0] <= rest_smallest then inputs[0] else rest_smallest;
    assert inputs == [inputs[0]] + inputs[1..]; 
    assert ret == if inputs[0] <= rest_smallest then inputs[0] else rest_smallest;
    ret 
}

function sum_array(inputs: seq<int>): int
{
  if |inputs| == 0 then 0
  else  inputs[0] + sum_array(inputs[1..])
}
\end{lstlisting}
\clearpage
\section*{GenAI interactions}
Below you will find all interactions between the `user' (program) and the `assistant' (OpenAI).
\subsection*{Program $\rightarrow$ GenAI}
\begin{lstlisting}
You are given the following task to perform in Dafny:

Write a method that computes the sum of an array of values, except for the smallest one, in a single loop. In the loop, update the sum and the smallest value. After the loop, return the difference.

The signature should be:

p_6_4_sum_without_smallest_array(inputs: array<int>) returns (sum: int)

The method should respect the following contract:

requires inputs.Length >= 1, ensures sum == sum_array(inputs[..]) - find_smallest(inputs[..])

The contract uses the following dafny code:

function find_smallest(inputs: seq<int>): (ret: int)
  requires |inputs| >= 1
  ensures forall x :: x in inputs ==> ret <= x 
{
  if |inputs| == 1 then inputs[0]
  else 
    assert 2 <= |inputs|;
    var rest_smallest := find_smallest(inputs[1..]);
    var ret: int := if inputs[0] <= rest_smallest then inputs[0] else rest_smallest;
    assert inputs == [inputs[0]] + inputs[1..]; 
    assert ret == if inputs[0] <= rest_smallest then inputs[0] else rest_smallest;
    ret 
}

function sum_array(inputs: seq<int>): int
{
  if |inputs| == 0 then 0
  else  inputs[0] + sum_array(inputs[1..])
}

These function(s) must not be used in your implementation of the method.

Produce and show only the Dafny body of this method, including the curly braces that surround it. Do not show the signature nor contract. You must not use 'assume' nor 'expect' anywhere in your code.
\end{lstlisting}
\subsection*{GenAI $\rightarrow$ Program}
\textbf{System fingerprint: }null
\\\textbf{ID: }chatcmpl-Dn1XsphJxUqBjD7MDG4vaxaUzckm6
\\\textbf{Model: }gpt-5.5-2026-04-23
\\\textbf{Created at: }1780576384
\\\textbf{Finish reason: }stop
\\\textbf{Usage: }762 tokens in, and 4064 tokens out
\begin{lstlisting}
{
  var i := inputs.Length - 1;
  var total := inputs[i];
  var smallest := inputs[i];

  assert 0 <= i < inputs.Length;
  assert |inputs[i..]| == 1;
  assert inputs[i..][0] == inputs[i];
  assert sum_array(inputs[i..]) == inputs[i];
  assert find_smallest(inputs[i..]) == inputs[i];

  while i > 0
    invariant 0 <= i < inputs.Length
    invariant total == sum_array(inputs[i..])
    invariant smallest == find_smallest(inputs[i..])
  {
    var j := i - 1;
    assert 0 <= j < inputs.Length;
    assert i == j + 1;
    assert |inputs[j..]| >= 2;
    assert inputs[j..][0] == inputs[j];
    assert inputs[j..][1..] == inputs[i..];

    assert sum_array(inputs[j..]) == inputs[j] + sum_array(inputs[i..]);
    assert find_smallest(inputs[j..]) ==
      (if inputs[j] <= find_smallest(inputs[i..]) then inputs[j] else find_smallest(inputs[i..]));

    total := inputs[j] + total;
    assert total == sum_array(inputs[j..]);

    if inputs[j] <= smallest {
      assert inputs[j] <= find_smallest(inputs[i..]);
      smallest := inputs[j];
    } else {
      assert !(inputs[j] <= find_smallest(inputs[i..]));
    }
    assert smallest == find_smallest(inputs[j..]);

    i := j;
  }

  assert i == 0;
  assert inputs[i..] == inputs[..];
  assert total == sum_array(inputs[..]);
  assert smallest == find_smallest(inputs[..]);
  sum := total - smallest;
}
\end{lstlisting}
\section*{Final Program}
\begin{lstlisting}
function find_smallest(inputs: seq<int>): (ret: int)
  requires |inputs| >= 1
  ensures forall x :: x in inputs ==> ret <= x 
{
  if |inputs| == 1 then inputs[0]
  else 
    assert 2 <= |inputs|;
    var rest_smallest := find_smallest(inputs[1..]);
    var ret: int := if inputs[0] <= rest_smallest then inputs[0] else rest_smallest;
    assert inputs == [inputs[0]] + inputs[1..]; 
    assert ret == if inputs[0] <= rest_smallest then inputs[0] else rest_smallest;
    ret 
}

function sum_array(inputs: seq<int>): int
{
  if |inputs| == 0 then 0
  else  inputs[0] + sum_array(inputs[1..])
}

method p_6_4_sum_without_smallest_array(inputs: array<int>) returns (sum: int)
	requires inputs.Length >= 1
	ensures sum == sum_array(inputs[..]) - find_smallest(inputs[..])
{
  var i := inputs.Length - 1;
  var total := inputs[i];
  var smallest := inputs[i];

  assert 0 <= i < inputs.Length;
  assert |inputs[i..]| == 1;
  assert inputs[i..][0] == inputs[i];
  assert sum_array(inputs[i..]) == inputs[i];
  assert find_smallest(inputs[i..]) == inputs[i];

  while i > 0
    invariant 0 <= i < inputs.Length
    invariant total == sum_array(inputs[i..])
    invariant smallest == find_smallest(inputs[i..])
  {
    var j := i - 1;
    assert 0 <= j < inputs.Length;
    assert i == j + 1;
    assert |inputs[j..]| >= 2;
    assert inputs[j..][0] == inputs[j];
    assert inputs[j..][1..] == inputs[i..];

    assert sum_array(inputs[j..]) == inputs[j] + sum_array(inputs[i..]);
    assert find_smallest(inputs[j..]) ==
      (if inputs[j] <= find_smallest(inputs[i..]) then inputs[j] else find_smallest(inputs[i..]));

    total := inputs[j] + total;
    assert total == sum_array(inputs[j..]);

    if inputs[j] <= smallest {
      assert inputs[j] <= find_smallest(inputs[i..]);
      smallest := inputs[j];
    } else {
      assert !(inputs[j] <= find_smallest(inputs[i..]));
    }
    assert smallest == find_smallest(inputs[j..]);

    i := j;
  }

  assert i == 0;
  assert inputs[i..] == inputs[..];
  assert total == sum_array(inputs[..]);
  assert smallest == find_smallest(inputs[..]);
  sum := total - smallest;
}
\end{lstlisting}
\section*{Total Token Usage}
\textbf{Input tokens: }762
\\\textbf{Output tokens: }4064
\\\textbf{Reasoning tokens: }3574
\\\textbf{Sum of `total tokens': }4826
\section*{Experiment Timings}
\textbf{Overall Experiment} started at 1780576383607, ended at 1780576475976, lasting 92369ms (92.37 seconds)
\\\textbf{Iteration \#1} started at 1780576383622, ended at 1780576475976, lasting 92354ms (92.35 seconds)
\end{document}
